\documentclass[12pt]{article}
\usepackage[utf8]{inputenc}
\usepackage[T1]{fontenc}
\usepackage{hyperref}
\title{Sindang Kelingi Language Resources for K-12 Curriculum Development}
\author{Compiled by Claude Sonnet 5.5 (claude-sonnet-5-5) for Andr\'as Kornai}
\date{}
\begin{document}
\maketitle

\section{Introduction}
Sindang Kelingi (ISO 639-3 sdi; Glottocode sind1280) is a Malayic variety of southern Sumatra, spoken in Musi Rawas Regency, Lubuklinggau City and neighbouring parts of Rejang Lebong and Empat Lawang, along the Kelingi, Beliti and Lakitan rivers. The code sdi was retired and merged into Col [liw] in 2008. Local usage calls the language Sindang, or Col/Cul after a word meaning ``no''. It is distinct from Rejang and Lampung. The prompt's family label (``Bookkeeping'') appears to be a data error. Glottolog's page title uses the same word, which suggests the error comes from that source.

\section{Primary Educational Resources}
No primers, textbooks, dictionaries or local-content (muatan lokal) materials were found. A search for them returned only academic articles.

\section{Linguistic and Reference Materials}
\begin{itemize}
\item Glottolog entry for Sindang Kelingi: \url{https://glottolog.org/resource/languoid/id/sind1280}
\item ``Bahasa Sindang: Studi Historis dan Kearifan Lokal'' (historical and cultural study of Sindang): \url{https://jurnal.radenfatah.ac.id/index.php/Suluh/article/view/21059}
\item ``Struktur Kalimat Tunggal Bahasa Sindang di Kota Lubuklinggau dan Pengaruhnya dalam Pembelajaran Bahasa Indonesia'' (simple sentence structure in Lubuklinggau Sindang): \url{https://oalib-perpustakaan.upi.edu/Record/doaj_61785753aba449edb116bdf92c52b203}
\item Same journal article, full text (AKSIS): \url{https://journal.unj.ac.id/unj/index.php/aksis/article/download/3064/2211}
\item ``Fungsi Bahasa dalam Acara Berasan Masyarakat Sindang Kelingi Musi Rawas'' (language in the berasan ceremony): \url{https://journal.ipm2kpe.or.id/index.php/KIBASP/article/view/908}
\item History of Lubuklinggau under the Sindang Kelingi Ilir clan, 1855--1942 (historical context): \url{https://ejournal.unibabwi.ac.id/index.php/santhet/article/view/5564}
\end{itemize}

\section{Religious and Community Resources}
\begin{itemize}
\item Joshua Project profile of the people group (listed under Col): \url{https://joshuaproject.net/people_groups/14191/ID}
\item Global Recordings Network lists a ``Musi: Kelingi'' language page: \url{https://www.globalrecordings.net/en/language/25907}
\end{itemize}
No scripture translation was confirmed.

\section{Digital and Audio Resources}
The Global Recordings Network page above is the only audio lead. Its content was not checked.

\section{Curriculum Development Context}
Indonesian is the language of schooling. Local-language teaching in Indonesia is done through local-content subjects. No Sindang materials were found. Leads for the next step are the education offices of Lubuklinggau and Musi Rawas, and searches under the names Sindang, Col/Cul and Lembak. Sources disagree on naming, so the Col [liw] and Sindang literatures should be checked together.
\end{document}